\documentclass[11pt,a4paper]{article}
\usepackage[T1]{fontenc}
\usepackage{lmodern,amsmath,amssymb,amsthm,mathtools}
\usepackage[margin=26mm]{geometry}
\usepackage{microtype,booktabs,array,xcolor,listings,hyperref}
\newcolumntype{P}[1]{>{\raggedright\arraybackslash}p{#1}}
\definecolor{navy}{RGB}{25,48,69}
\hypersetup{colorlinks=true,linkcolor=navy,urlcolor=navy,citecolor=navy,
  pdftitle={Parameter Monotonicity of Mittag-Leffler Functions: the Conjecture on (2,4) and the Sharp Range},
  pdfauthor={Prepared for Keio}}
\theoremstyle{plain}
\newtheorem{theorem}{Theorem}[section]
\newtheorem{lemma}[theorem]{Lemma}
\newtheorem{corollary}[theorem]{Corollary}
\theoremstyle{remark}
\newtheorem*{remark}{Remark}
\newcommand{\E}{\mathbb E}
\newcommand{\R}{\mathbb R}
\newcommand{\Z}{\mathbb Z}
\newcommand{\C}{\mathbb C}
\DeclareMathOperator{\Log}{Log}
\DeclareMathOperator{\Res}{Res}
\lstset{language=Python,basicstyle=\ttfamily\scriptsize,columns=fullflexible,keepspaces=true,
  breaklines=true,showstringspaces=false,frame=single,framerule=0.3pt,rulecolor=\color{gray},
  upquote=true,tabsize=4,xleftmargin=2pt,xrightmargin=2pt}
\setlength{\emergencystretch}{2em}
\allowdisplaybreaks

\title{\color{navy}Parameter Monotonicity of Mittag--Leffler Functions:\\
the Conjecture on $(2,4)$ and the Sharp Range $(0,\alpha^*]$}
\author{Prepared for Keio\thanks{Section~\ref{sec:A} reproduces the computer-assisted draft
\emph{Parameter monotonicity of Mittag--Leffler functions on $(0,4]$} (14 November 2026); its
certificate was re-implemented independently for this version. Sections~\ref{sec:B}
and~\ref{sec:sharp} and the unified presentation were prepared with AI assistance
(Claude, Anthropic) on 15 November 2026. No external review has taken place.}}
\date{15 November 2026\\\small Unified computer-assisted proof draft for independent review}

\begin{document}
\maketitle

\begin{abstract}
For $x>0$ put $U_\alpha(x)=E_\alpha(x^\alpha)$, where $E_\alpha$ is the Mittag--Leffler
function. Ferreira and Simon proved that $\alpha\mapsto U_\alpha(x)$ decreases on $(0,2)$,
expected the decrease to persist on $(2,4)$, and showed that it fails for some large $\alpha$.
We prove the conjecture on $[2,4]$ twice. Both proofs start from the same Laplace kernel
$K_\alpha(s)=s^{\alpha-1}\Log s/(s^\alpha-1)^2$ and differ in where its Bromwich contour is moved.
Proof~A shifts the line to $\Re s=1/2$ and certifies a compact rectangle. Proof~B rotates the
contour onto two rays in the left half-plane, keeps the residues at $e^{\pm2\pi i/\alpha}$, and
needs only a one-dimensional certificate. Combining the two techniques, we determine the maximal
range: for every $x>0$ the map $\alpha\mapsto U_\alpha(x)$ is strictly decreasing on
$(0,8.8033]$, while $\partial_\alpha U_\alpha(x)>0$ at $(\alpha,x)=(8.8034,10.7539)$. Hence the
threshold $\alpha^*$ satisfies $8.8033\le\alpha^*<8.8034$; numerically
$\alpha^*=8.80337769817\ldots$. Consequences include strict monotonicity of
$E_{\alpha,\beta}(x^\alpha)$ for every $\beta\ge1$ on $(0,8.8033]$ and two-sided bounds by
elementary functions. All finite verifications use Arb ball arithmetic, and the programs are
reproduced in the appendix. No claim of external acceptance or of historical priority is made.
\end{abstract}

\tableofcontents

\section{Introduction}\label{sec:intro}
\subsection{Background}
For $\alpha,\beta>0$ let
\[
 E_{\alpha,\beta}(z)=\sum_{n=0}^{\infty}\frac{z^n}{\Gamma(\beta+\alpha n)},\qquad
 E_\alpha=E_{\alpha,1},\qquad U_\alpha(x)=E_\alpha(x^\alpha)\quad(x>0).
\]
The scaling $x^\alpha$ puts all orders on one exponential scale: $U_\alpha(x)\sim e^x/\alpha$
as $x\to\infty$, and
\[
 U_1(x)=e^x,\quad U_2(x)=\cosh x,\quad
 U_3(x)=\tfrac13\bigl(e^x+2e^{-x/2}\cos\tfrac{\sqrt3\,x}{2}\bigr),\quad
 U_4(x)=\tfrac12(\cosh x+\cos x).
\]
Equivalently, $U_\alpha$ solves the Caputo problem $D^\alpha u=u$, $u(0)=1$,
$u^{(j)}(0)=0$ for $1\le j<\lceil\alpha\rceil$ \cite{GKMR}. Monotonicity in $\alpha$ therefore
compares solutions of fractional growth equations of different orders.

Ferreira and Simon \cite[Theorem~A]{FS} proved that $\alpha\mapsto U_\alpha(x)$ decreases on
$(0,2)$ for every $x>0$, through convex ordering of an associated family of random variables.
In \cite[Remark~1(a)]{FS} they discuss larger orders. They expect, supported by simulations, that
the decrease persists on $(2,4)$, and they check $U_4<U_3<U_2$ from the closed forms above. They
also explain why their probabilistic argument does not reach $\alpha>2$, and observe numerically
that $U_{25}(30)<U_{26}(30)$, so monotonicity fails for some large $\alpha$. Two questions remain:
the conjecture on $(2,4)$, and the exact range of $\alpha$ on which the decrease holds for all $x>0$.

\subsection{Main results}
Write $H_\alpha(x)=-\partial_\alpha U_\alpha(x)$ and $\psi=\Gamma'/\Gamma$.

\begin{theorem}\label{thm:24}
For every $x>0$ and every $\alpha\in[2,4]$, $H_\alpha(x)>0$. In particular
$\alpha\mapsto E_\alpha(x^\alpha)$ is strictly decreasing on $[2,4]$, and the conjecture of
\cite[Remark~1(a)]{FS} holds.
\end{theorem}

\begin{theorem}\label{thm:sharp}
\begin{enumerate}
 \item[(i)] For every $x>0$, the map $\alpha\mapsto E_\alpha(x^\alpha)$ is strictly decreasing on
 $(0,8.8033]$.
 \item[(ii)] $\partial_\alpha E_\alpha(x^\alpha)>0$ at $(\alpha,x)=(8.8034,\,10.7539)$.
\end{enumerate}
Consequently the threshold
\[
 \alpha^*:=\sup\bigl\{A>0:\ \alpha\mapsto E_\alpha(x^\alpha)\ \text{is non-increasing on }(0,A]
 \text{ for every }x>0\bigr\}
\]
satisfies $8.8033\le\alpha^*<8.8034$. Numerically $\alpha^*=8.80337769817\ldots$, attained at
$x^*=10.7538996\ldots$ (Section~\ref{subsec:critical}).
\end{theorem}

\begin{corollary}\label{cor:beta}
For every $\beta\ge1$ and $x>0$, the map $\alpha\mapsto E_{\alpha,\beta}(x^\alpha)$ is strictly
decreasing on $(0,8.8033]$.
\end{corollary}

\begin{corollary}\label{cor:bounds}
For every $x>0$,
\begin{gather*}
 \frac{\cosh x+\cos x}{2}<E_\alpha(x^\alpha)<\cosh x\qquad(2<\alpha<4),\\
 E_8(x^8)<E_\alpha(x^\alpha)<\frac{\cosh x+\cos x}{2}\qquad(4<\alpha<8),
\end{gather*}
where
$E_8(x^8)=\tfrac14\bigl(\cosh x+\cos x+2\cos\tfrac{x}{\sqrt2}\cosh\tfrac{x}{\sqrt2}\bigr)$.
\end{corollary}

\subsection{Structure of the proofs}
Every argument below uses the Laplace transform of $H_\alpha$,
\[
 K_\alpha(s)=\frac{s^{\alpha-1}\Log s}{(s^\alpha-1)^2}\qquad(\Re s>1),
\]
and its Bromwich inversion (Lemmas~\ref{lem:laplace} and~\ref{lem:inversion}); the proofs differ
in where the contour is moved.
\begin{itemize}
 \item Proof~A (Section~\ref{sec:A}) shifts the vertical line to $\Re s=1/2$. This crosses only
 the pole at $s=1$, but the remainder is of size $e^{x/2}$, so a compact rectangle in $(\alpha,x)$
 is certified.
 \item Proof~B (Section~\ref{sec:B}) rotates the contour onto two rays in the left half-plane and
 keeps the residues at $e^{\pm2\pi i/\alpha}$ explicitly. The remainder then stays bounded in $x$,
 and the certified part is one-dimensional.
 \item For $\alpha>4$ more residues enter (Section~\ref{sec:sharp}). The rays of Proof~B are then
 combined with a two-dimensional certificate in the style of Proof~A, which is needed only near the
 critical point.
\end{itemize}
Small $x$ is handled in every case by termwise positivity of the series for $H_\alpha$
(Lemma~\ref{lem:small}). Table~\ref{tab:cert} lists the finite certified computations.

\begin{table}[ht]\centering\small\setlength{\tabcolsep}{4pt}
\begin{tabular}{llll}
\toprule
Step & Region & Certificate & Certified margin\\
\midrule
Lemma~\ref{lem:certA} (Proof A) & $[2,4]\times[\frac52,16]$ & 2-D, 16{,}749 rectangles &
 $e^{-x}H_{\alpha,24}>10^{-5}$\\
Lemma~\ref{lem:certB} (Proof B) & $\alpha\in[2,4]$ & 1-D, 2{,}000 intervals & $G\le0.1999$\\
Lemma~\ref{lem:tail16} & $\alpha\in[4,8.8033]$, $x\ge16$ & 1-D, 960 intervals & $G\le0.5005$\\
Lemma~\ref{lem:middle} & $\alpha\in[4,8.8033]$, $e^{\psi(1+\alpha)}\le x\le16$ & 2-D, 385 boxes &
 $Q\ge1.21\cdot10^{-5}$\\
Lemma~\ref{lem:witness} & $(\alpha,x)=(8.8034,10.7539)$ & single point &
 $\partial_\alpha U_\alpha\ge0.00396$\\
\bottomrule
\end{tabular}
\caption{Finite certified computations. All use Arb ball arithmetic \cite{Arb} through
python-flint~0.9.0 \cite{flint}; the programs are listed in Appendix~\ref{app:code}. Here
$G<1$ and $Q>0$ are the sufficient conditions of Lemma~\ref{lem:bounds}.}\label{tab:cert}
\end{table}

\subsection{Status}
This is a draft for independent mathematical and software review. The certified steps are
ball-arithmetic evaluations of elementary functions and of $\Gamma$, $\log\Gamma$ and $\psi$; they
are not formalized in a proof assistant. The literature check was limited: as of
15 November 2026 we located \cite{FS} (version~1, without journal reference) and no published
resolution of the questions above. This does not establish priority.

\section{Analytic preliminaries}\label{sec:prelim}
Throughout, $\Log s$ and $s^\beta=e^{\beta\Log s}$ are principal branches on
$\C\setminus(-\infty,0]$.

\begin{lemma}[derivative series; small $x$]\label{lem:small}
For $x>0$ and $\alpha>0$,
\begin{equation}\label{eq:series}
 H_\alpha(x)=\sum_{n=1}^{\infty}\frac{n\,x^{\alpha n}}{\Gamma(1+\alpha n)}
 \bigl(\psi(1+\alpha n)-\log x\bigr),
\end{equation}
with locally uniform convergence in $(\alpha,x)$. If $0<x\le e^{\psi(1+\alpha)}$, every term is
nonnegative and every term with $n\ge2$ is positive; hence $H_\alpha(x)>0$.
\end{lemma}
\begin{proof}
The series for $U_\alpha$ and its termwise $\alpha$-derivative converge locally uniformly because
$\Gamma(1+\alpha n)$ grows faster than any exponential in $n$, and
$\partial_\alpha\bigl[x^{\alpha n}/\Gamma(1+\alpha n)\bigr]
 =n\,x^{\alpha n}\bigl(\log x-\psi(1+\alpha n)\bigr)/\Gamma(1+\alpha n)$.
Since $\psi'(y)=\sum_{j\ge0}(y+j)^{-2}>0$, we have
$\psi(1+\alpha n)\ge\psi(1+\alpha)\ge\log x$ for $n\ge1$, with strict first inequality for $n\ge2$.
\end{proof}

For $\alpha\in[2,8.8033]$ the threshold $e^{\psi(1+\alpha)}$ increases from
$e^{\psi(3)}=2.5162\ldots$ to $e^{\psi(9.8033)}=9.3077\ldots$; at $\alpha=4$ it equals
$e^{\psi(5)}=4.5091\ldots$.

\begin{lemma}[Laplace transform]\label{lem:laplace}
For $\alpha\ge2$ and $\Re s>1$, $\displaystyle\int_0^\infty e^{-sx}H_\alpha(x)\,dx=K_\alpha(s)$.
\end{lemma}
\begin{proof}
Let $c=\Re s>1$, $n\ge1$ and $y=1+\alpha n\ge1+2n$. If $Y$ has the gamma law of shape $y$ and unit
rate, the substitution $x=u/c$ gives
\begin{align*}
 \int_0^\infty e^{-cx}\frac{n\,x^{\alpha n}}{\Gamma(1+\alpha n)}\bigl|\psi(y)-\log x\bigr|\,dx
 &=n\,c^{-y}\,\E\bigl|\psi(y)-\log Y+\log c\bigr|\\
 &\le n\,c^{-1-2n}\bigl(\psi(y)+y+(y-1)^{-1}+\log c\bigr),
\end{align*}
using $0<\psi(y)=\E\log Y$ and $\E|\log Y|\le\E(Y+Y^{-1})=y+(y-1)^{-1}$. The bound is
$O(n^2c^{-2n})$ locally uniformly in $\alpha$, hence summable, and termwise integration is
justified. Differentiating $\int_0^\infty e^{-sx}x^\nu\,dx=\Gamma(\nu+1)s^{-\nu-1}$ in $\nu$ gives
$\int_0^\infty e^{-sx}x^\nu\bigl(\psi(\nu+1)-\log x\bigr)dx/\Gamma(\nu+1)=s^{-\nu-1}\Log s$.
With $\nu=\alpha n$ and $|s^{-\alpha}|<1$,
$\sum_{n\ge1}n\,s^{-1-\alpha n}\Log s=s^{\alpha-1}\Log s/(s^\alpha-1)^2$.
\end{proof}

\begin{lemma}[inversion]\label{lem:inversion}
For $\alpha\ge2$, $c>1$ and $x>0$,
\[
 H_\alpha(x)=\frac{1}{2\pi i}\int_{c-i\infty}^{c+i\infty}e^{xs}K_\alpha(s)\,ds,
\]
and the integral converges absolutely.
\end{lemma}
\begin{proof}
On $\Re s=c$ we have $|K_\alpha(c+it)|=O(|t|^{-\alpha-1}\log|t|)$, so $t\mapsto K_\alpha(c+it)$ is
integrable. By Lemma~\ref{lem:laplace} it is the Fourier transform of
$f(x)=e^{-cx}H_\alpha(x)\mathbf 1_{x>0}\in L^1(\R)$. Fourier inversion recovers $f$ almost
everywhere; both sides are continuous on $(0,\infty)$.
\end{proof}

\begin{lemma}[poles and residues]\label{lem:poles}
On $\C\setminus(-\infty,0]$, $K_\alpha$ is meromorphic, with poles exactly at $w_k=e^{iu_k}$,
$u_k=2\pi k/\alpha$, $k\in\Z$, $|u_k|<\pi$. The pole $w_0=1$ is simple, the others are double, and
\[
 \Res_{s=w_k}e^{xs}K_\alpha(s)=\frac{e^{xw_k}}{\alpha^2}\bigl(1+iu_k\,x\,w_k\bigr).
\]
\end{lemma}
\begin{proof}
$s^\alpha=1$ if and only if $\alpha\Log s\in2\pi i\Z$, and these zeros are simple because
$(s^\alpha)'=\alpha s^{\alpha-1}\ne0$. Put $g(s)=e^{xs}s^{\alpha-1}\Log s$ and $F(s)=s^\alpha-1$.
At a simple zero $p$ of $F$, $\Res(g/F^2)=g'(p)/F'(p)^2-g(p)F''(p)/F'(p)^3$. With $p^\alpha=1$,
$F'(p)=\alpha/p$, $F''(p)=\alpha(\alpha-1)/p^2$ and $\Log p=iu_k$, this equals
$e^{xp}(1+x\,p\Log p)/\alpha^2$. At $p=1$ the zero of $\Log s$ cancels one order and the same
formula gives $e^x/\alpha^2$.
\end{proof}

\begin{lemma}[rotated rays]\label{lem:rays}
Let $\alpha\ge2$, $x>0$ and $\varphi\in(\pi/2,\pi)$ with $\alpha\varphi\notin2\pi\Z$, and put
$K_\varphi=\#\{k\ge1:u_k<\varphi\}$. Then
\begin{equation}\label{eq:rays}
 H_\alpha(x)=\frac{e^x+S_\varphi(\alpha,x)}{\alpha^2}+J_\varphi(\alpha,x),
\end{equation}
where
\[
 S_\varphi(\alpha,x)=\sum_{k=1}^{K_\varphi}2e^{x\cos u_k}
 \Bigl[\cos(x\sin u_k)-u_kx\,\sin(x\sin u_k+u_k)\Bigr],
\]
\[
 J_\varphi(\alpha,x)=\frac1\pi\,\Im\int_0^\infty e^{xre^{i\varphi}}\,r^{\alpha-1}e^{i\alpha\varphi}\,
 \frac{\log r+i\varphi}{\bigl(r^\alpha e^{i\alpha\varphi}-1\bigr)^2}\,dr .
\]
\end{lemma}
\begin{proof}
Let $\Gamma_\varphi$ run from $\infty e^{-i\varphi}$ to $0$ and on to $\infty e^{i\varphi}$, and
consider the region $\{|\arg s|<\varphi,\ \Re s<c\}$ between $\Gamma_\varphi$ and the line of
Lemma~\ref{lem:inversion}; the whole line lies in the sector because $|\arg(c+it)|<\pi/2$.
$K_\alpha$ has no pole on $\Gamma_\varphi$, since $s^\alpha=1$ on $\arg s=\pm\varphi$ would force
$\alpha\varphi\in2\pi\Z$. On arcs $|s|=R$ with $\Re s\le c$ we have $|e^{xs}|\le e^{cx}$ and
$|K_\alpha(s)|=O(R^{-\alpha-1}\log R)$, so these arcs contribute nothing as $R\to\infty$. Near
$s=0$, $|K_\alpha(s)|=O(|s|^{\alpha-1}|\log|s||)$, so a small arc around the origin also
contributes nothing. Every pole has modulus $1<c$, so the residue theorem gives
\[
 H_\alpha(x)=\sum_{|u_k|<\varphi}\Res_{s=w_k}e^{xs}K_\alpha(s)
 +\frac1{2\pi i}\int_{\Gamma_\varphi}e^{xs}K_\alpha(s)\,ds .
\]
By Lemma~\ref{lem:poles} the residues at $w_{\pm k}$ are complex conjugate, and their sum is the
$k$-th term of $S_\varphi/\alpha^2$. Since $K_\alpha(\bar s)=\overline{K_\alpha(s)}$, the ray
integral equals $J_\varphi$. The argument stays inside the open sector $|\arg s|<\varphi<\pi$ and
never meets the branch cut, so it holds for even integer $\alpha$ as well.
\end{proof}

\begin{lemma}[explicit bounds]\label{lem:bounds}
In the setting of Lemma~\ref{lem:rays} put $\lambda=-x\cos\varphi>0$ and $\gamma=\alpha\varphi$,
and let $m^2=1$ if $\cos\gamma\le0$ and $m^2=\sin^2\gamma$ otherwise. Then
$|J_\varphi(\alpha,x)|\le B_\varphi(\alpha,x)$, where
\[
 B_\varphi(\alpha,x)=\frac{1}{\pi m^2}\Bigl[\frac{\Gamma(\alpha)}{\lambda^\alpha}
 \bigl(\varphi+(\psi(\alpha)-\log\lambda)_+\bigr)
 +2\min\Bigl(\frac1{\alpha^2},\frac{\Gamma(\alpha-1)}{e\,\lambda^{\alpha-1}}\Bigr)\Bigr],
\]
and $|e^{-x}S_\varphi(\alpha,x)|\le\sum_{k=1}^{K_\varphi}2e^{-x(1-\cos u_k)}(1+u_kx)$.
Consequently
\[
 \alpha^2e^{-x}H_\alpha(x)\ \ge\ Q_\varphi(\alpha,x):=1+e^{-x}S_\varphi(\alpha,x)
 -\alpha^2e^{-x}B_\varphi(\alpha,x)\ \ge\ 1-G_\varphi(\alpha,x),
\]
\[
 G_\varphi(\alpha,x):=\sum_{k=1}^{K_\varphi}2e^{-x(1-\cos u_k)}(1+u_kx)
 +\alpha^2e^{-x}B_\varphi(\alpha,x).
\]
In particular $H_\alpha(x)>0$ whenever $Q_\varphi(\alpha,x)>0$, and in particular whenever
$G_\varphi(\alpha,x)<1$.
\end{lemma}
\begin{proof}
On $s=re^{\pm i\varphi}$ we have $|e^{xs}|=e^{-\lambda r}$, $|s^{\alpha-1}|=r^{\alpha-1}$,
$|\log r\pm i\varphi|\le|\log r|+\varphi$ and
$|r^\alpha e^{\pm i\gamma}-1|^2=(r^\alpha-\cos\gamma)^2+\sin^2\gamma\ge m^2$. Hence
$|J_\varphi|\le(\pi m^2)^{-1}\int_0^\infty e^{-\lambda r}r^{\alpha-1}(|\log r|+\varphi)\,dr$. Write
$|\log r|=\log r+2\log(1/r)\mathbf 1_{r<1}$ and use
$\int_0^\infty e^{-\lambda r}r^{\alpha-1}\log r\,dr=\Gamma(\alpha)\lambda^{-\alpha}(\psi(\alpha)-\log\lambda)$.
The remaining integral $L=\int_0^1e^{-\lambda r}r^{\alpha-1}\log(1/r)\,dr$ is at most
$\int_0^1r^{\alpha-1}\log(1/r)\,dr=\alpha^{-2}$, and, since $\log(1/r)\le 1/(er)$, at most
$\Gamma(\alpha-1)/(e\lambda^{\alpha-1})$. Finally
$|e^{xw}(1+iuxw)|\le e^{x\cos u}(1+ux)$ for $w=e^{iu}$.
\end{proof}

\begin{lemma}[monotonicity in $x$]\label{lem:mono}
Let $2\le\alpha\le8.8034$ and let $\varphi$ be as in Lemma~\ref{lem:rays}. Then
$x\mapsto G_\varphi(\alpha,x)$ is strictly decreasing on $[4.1,\infty)$, and on $[1,\infty)$ if
$\alpha\le4$.
\end{lemma}
\begin{proof}
For $c=\cos u$, the derivative of $e^{-x(1-c)}(1+ux)$ is $e^{-x(1-c)}\bigl(u-(1-c)(1+ux)\bigr)$,
which is negative for $x>(1-c)^{-1}-u^{-1}$. Here $u\ge u_1=2\pi/\alpha\ge0.7137$, so
$1-c\ge1-\cos0.7137>0.244$ and $(1-c)^{-1}<4.1$; if $\alpha\le4$ then $c\le0$ and $(1-c)^{-1}\le1$.
In $B_\varphi$, $\lambda$ is proportional to $x$, and the functions $\lambda^{-\alpha}$,
$\lambda^{-\alpha}(\psi(\alpha)-\log\lambda)_+$ (derivative
$-\lambda^{-\alpha-1}[1+\alpha(\psi(\alpha)-\log\lambda)]<0$ where it is positive) and the minimum are
nonincreasing in $\lambda$. The factor $e^{-x}$ makes $e^{-x}B_\varphi$ strictly decreasing.
\end{proof}

\section{Proof A of Theorem~\ref{thm:24}: the line \texorpdfstring{$\Re s=1/2$}{Re s = 1/2}}\label{sec:A}
This section follows the draft of 14 November 2026. Its compact certificate was re-implemented
independently (Appendix~\ref{app:A}) and reproduces the draft's data exactly.

\begin{lemma}\label{lem:line}
Uniformly for $\alpha\in[2,4]$ and $x>0$, $\bigl|H_\alpha(x)-e^x/\alpha^2\bigr|<100\,e^{x/2}$.
In particular $H_\alpha(x)>0$ for $x\ge16$.
\end{lemma}
\begin{proof}
For $\alpha\in[2,4]$ the only pole of $K_\alpha$ in $\Re s>0$ is $s=1$: there
$|\arg s|<\pi/2$, so $\alpha\arg s$ cannot be a nonzero multiple of $2\pi$. At $\alpha=4$ the
further roots $\pm i$ lie to the left of both lines below. Moving the line of
Lemma~\ref{lem:inversion} to $\Re s=1/2$ crosses only $s=1$, and the horizontal segments vanish by
the decay of $K_\alpha$. Hence
\[
 H_\alpha(x)=\frac{e^x}{\alpha^2}+\frac1{2\pi}\int_{-\infty}^{\infty}
 e^{(1/2+it)x}K_\alpha\bigl(\tfrac12+it\bigr)\,dt .
\]
Let $s=\frac12+it$, $r=|s|$, and take $t\ge0$ (conjugation handles $t<0$).
\begin{enumerate}
 \item If $0\le t\le\frac12$: $r\le2^{-1/2}$, $|s^\alpha-1|\ge1-r^\alpha\ge\frac12$,
 $|\Log s|\le\log2+\frac\pi4<\frac32$ and $r^{\alpha-1}\le1$, so $|K_\alpha(s)|<6$.
 \item If $\frac12\le t\le2$: $\frac\pi4\le\arg s\le\arctan4$, and
 $\frac\pi2\le\alpha\arg s<\frac{11\pi}6$ because $\arctan\frac14\ge\frac4{17}>\frac\pi{24}$. If
 $\cos(\alpha\arg s)\le0$ the distance from $s^\alpha$ to $1$ is at least $1$; otherwise
 $\alpha\arg s\in(\frac{3\pi}2,\frac{11\pi}6)$ and $|\sin(\alpha\arg s)|\ge\frac12$. Thus
 $|s^\alpha-1|\ge\frac12$. With $r\le\sqrt{17}/2$, $r^{\alpha-1}\le\max(r,r^3)<9$ and
 $|\Log s|<3$ we get $|K_\alpha(s)|<108<144$.
 \item If $t\ge2$: $r^\alpha\ge4$, $|s^\alpha-1|\ge\frac34r^\alpha$, $\log(r/t)\le\frac1{8t^2}<\frac14$
 and $|\arg s|<\frac\pi2<2$, so $|K_\alpha(s)|<2t^{-3}(\log t+\frac94)$.
\end{enumerate}
Therefore
$\int_\R|K_\alpha(\frac12+it)|\,dt<2\bigl(3+216+\frac{4\log2+11}{16}\bigr)<440$, and
$|H_\alpha(x)-e^x/\alpha^2|<\frac{440}{2\pi}e^{x/2}<100e^{x/2}$. For $x\ge16$,
$H_\alpha(x)>e^x/16-100e^{x/2}>0$ because $e^8>1600$.
\end{proof}

\begin{lemma}\label{lem:scalar}
$\psi(3)>\log\frac52$ and $\psi(51)>\log16$. Consequently, for $\alpha\in[2,4]$, every term
of~\eqref{eq:series} is positive when $0<x\le\frac52$, and every term with $n\ge25$ is positive
when $x\le16$. In particular $H_\alpha(x)\ge H_{\alpha,24}(x)$ for $x\in[\frac52,16]$, where
\[
 H_{\alpha,24}(x)=\sum_{n=1}^{24}\frac{n\,x^{\alpha n}}{\Gamma(1+\alpha n)}
 \bigl(\psi(1+\alpha n)-\log x\bigr).
\]
\end{lemma}
\begin{proof}
Arb gives $\psi(3)-\log\frac52\in[0.00649360\pm3.3\cdot10^{-9}]$ and
$\psi(51)-\log16\in[1.14940\pm10^{-6}]$ (Appendix~\ref{app:A}); note $\psi(3)=\frac32-\gamma_E$.
The rest follows from $1+\alpha n\ge3$, from $1+\alpha n\ge51$ for $n\ge25$, and from the
monotonicity of $\psi$.
\end{proof}

\begin{lemma}[compact certificate]\label{lem:certA}
For every $(\alpha,x)\in[2,4]\times[\frac52,16]$, $e^{-x}H_{\alpha,24}(x)>10^{-5}$.
\end{lemma}
\begin{proof}
The program of Appendix~\ref{app:A} starts from the exact rational rectangle
$[2,4]\times[\frac52,16]$. On each rectangle it encloses
\[
 \sum_{n=1}^{24}n\exp\bigl(nA\log X-\log\Gamma(1+nA)-X\bigr)\bigl(\psi(1+nA)-\log X\bigr)
\]
in 96-bit Arb ball arithmetic, with $A$ and $X$ balls containing the parameter intervals, and
accepts the rectangle if the lower end of the ball exceeds the exact rational $10^{-5}$.
Otherwise it bisects $\alpha$ when eight times its width exceeds the $x$-width, and $x$ otherwise.
All endpoints are exact rationals, each split replaces a rectangle by two with the same union, and
a depth guard can only cause failure. The run terminates with an empty stack after
33{,}497 evaluations and accepts 16{,}749 rectangles of total exact area $27$, with maximum depth
$16$ (so $33497=2\cdot16749-1$). The smallest certified lower bound is $1.1489\cdot10^{-5}$.
\end{proof}

The draft reported the same rectangle data from its own program, together with a replay at 160
bits using direct powers and gamma values (34{,}246 subrectangles in 51{,}743 evaluations). That
replay is not part of this bundle.

\begin{proof}[Proof A of Theorem~\ref{thm:24}]
For $0<x\le\frac52$ use Lemma~\ref{lem:scalar}. For $x\in[\frac52,16]$, Lemmas~\ref{lem:scalar}
and~\ref{lem:certA} give $H_\alpha(x)\ge H_{\alpha,24}(x)>0$. For $x\ge16$ use
Lemma~\ref{lem:line}. These three ranges cover $x>0$ uniformly in $\alpha\in[2,4]$.
\end{proof}

\begin{remark}
The constant in Lemma~\ref{lem:line} is coarse. Numerically,
$\int_\R|K_\alpha(\frac12+it)|\,dt\le1.73$ for $\alpha\in[2,4]$, with the maximum at $\alpha=2$.
With this value and the $\alpha$-dependent threshold of Lemma~\ref{lem:small}, the ratio of the
remainder bound to $e^x/\alpha^2$ at $x=e^{\psi(1+\alpha)}$ is at most $0.31$. A certified version of
this one-dimensional integral bound would therefore remove the compact certificate from Proof~A.
This observation is not used.
\end{remark}

\section{Proof B of Theorem~\ref{thm:24}: rotated rays}\label{sec:B}
Use Lemma~\ref{lem:rays} with
\[
 \varphi=\tfrac{7\pi}{10}\ \text{on }[2,\tfrac52]\ (K_\varphi=0),\qquad
 \varphi=\tfrac{9\pi}{10}\ \text{on }[\tfrac52,4]\ (K_\varphi=1).
\]
On the first range $u_1=2\pi/\alpha\ge\frac{4\pi}5>\varphi$, so no residue besides $s=1$ enters. On
the second, $u_1\le\frac{4\pi}5<\varphi$, while $u_2=4\pi/\alpha\ge\pi$ is not a pole. In both
ranges $\alpha\varphi\notin2\pi\Z$, and $m^2\ge\frac12$ because $\gamma=\alpha\varphi$ lies in
$[1.4\pi,1.75\pi]$ and in $[2.25\pi,3.6\pi]$, respectively.

\begin{lemma}[one-dimensional certificate]\label{lem:certB}
Split $[2,\frac52]$ into $500$ and $[\frac52,4]$ into $1500$ equal intervals $[a,a']$. For each,
let $X$ be a certified lower bound of $e^{\psi(1+a)}$. Then
$\sup_{\alpha\in[a,a']}G_\varphi(\alpha,X)<1$. The largest certified upper bounds are $0.199825$
(on $[2.499,2.5]$) and $0.182407$ (on $[3.999,4]$).
\end{lemma}
\begin{proof}
The program of Appendix~\ref{app:B} evaluates $G_\varphi$ with $\alpha$ a ball containing $[a,a']$,
in 128-bit Arb arithmetic, with a certified lower bound of $m^2$. On every interval it also
asserts that $u_1$ lies on the stated side of $\varphi$ and that $X>1$.
\end{proof}

\begin{proof}[Proof B of Theorem~\ref{thm:24}]
Let $\alpha\in[2,4]$ and $x>0$. If $x\le e^{\psi(1+\alpha)}$, use Lemma~\ref{lem:small}. Otherwise,
for the interval $[a,a']\ni\alpha$ of Lemma~\ref{lem:certB}, monotonicity of $\psi$ gives
$x\ge e^{\psi(1+\alpha)}\ge X>2.5$. Lemmas~\ref{lem:mono} and~\ref{lem:certB} give
$G_\varphi(\alpha,x)\le G_\varphi(\alpha,X)<1$, and Lemma~\ref{lem:bounds} gives $H_\alpha(x)>0$.
\end{proof}

\begin{remark}
Proofs~A and~B share only Lemmas~\ref{lem:small}--\ref{lem:inversion}; their contour shifts,
bounds and certificates are independent. As a non-rigorous cross-check, formula~\eqref{eq:rays}
agrees with the series~\eqref{eq:series} to relative error at most $9.6\cdot10^{-41}$ in $36$ cases
($\alpha\in\{2.1,2.3,2.7,3,3.4,3.9\}$, $x\in\{0.7,2.5,6\}$, $\varphi\in\{0.75\pi,0.92\pi\}$,
40-digit arithmetic).
\end{remark}

\section{The sharp range: proof of Theorem~\ref{thm:sharp}}\label{sec:sharp}
By \cite[Theorem~A]{FS} and Theorem~\ref{thm:24}, it remains to treat $\alpha\in[4,8.8033]$ and to
certify the witness of part~(ii).

\subsection{Ray angles}
For $\alpha>4$ the poles $w_{\pm1}$ move into the right half-plane, and further poles $w_{\pm2}$,
$w_{\pm3}$, $w_{\pm4}$ appear on the principal sheet as $\alpha$ passes $4$, $6$ and $8$. We fix one
ray angle on each of five ranges (Table~\ref{tab:rays}).

\begin{table}[ht]\centering\small
\begin{tabular}{lcccc}
\toprule
$\alpha$-range & $\varphi/\pi$ & $K_\varphi$ & last pole inside, $u_{K_\varphi}/\pi$ &
 first pole outside, $u_{K_\varphi+1}/\pi$\\
\midrule
$[4,5]$ & $0.65$ & 1 & $[0.40,0.50]$ & $[0.80,1.00]$\\
$[5,6]$ & $0.90$ & 2 & $[0.667,0.80]$ & $[1.00,1.20]$\\
$[6,7]$ & $0.76$ & 2 & $[0.571,0.667]$ & $[0.857,1.00]$\\
$[7,8]$ & $0.66$ & 2 & $[0.50,0.571]$ & $[0.75,0.857]$\\
$[8,8.8033]$ & $0.83$ & 3 & $[0.682,0.75]$ & $[0.909,1.00]$\\
\bottomrule
\end{tabular}
\caption{Ray angles for $\alpha\in[4,8.8033]$; $u_k=2\pi k/\alpha$. A pole with $u_k\ge\pi$ is not
on the principal sheet.}\label{tab:rays}
\end{table}

\begin{lemma}[geometry]\label{lem:geom}
On each range of Table~\ref{tab:rays} and for every $\alpha$ in it: $\varphi\in(\pi/2,\pi)$,
$u_k<\varphi$ for $1\le k\le K_\varphi$, $u_{K_\varphi+1}>\varphi$, and $m^2>\frac12$. In
particular $\alpha\varphi\notin2\pi\Z$, and Lemma~\ref{lem:rays} applies with the stated
$K_\varphi$.
\end{lemma}
\begin{proof}
The inequalities for $u_k$ are checked in ball arithmetic on each whole range, and $m^2>\frac12$
on $400$ equal sub-intervals of each range (Appendix~\ref{app:C}). If $\alpha\varphi=2\pi k$ then
$\varphi=u_k$, which the inequalities exclude.
\end{proof}

\begin{lemma}[large $x$]\label{lem:tail16}
For every $\alpha\in[4,8.8033]$, $G_\varphi(\alpha,16)\le0.5005$. Hence $H_\alpha(x)>0$ for
$x\ge16$.
\end{lemma}
\begin{proof}
The bound is certified on $960$ equal sub-intervals: $200$ on each of the first four ranges and
$160$ on $[8,8.8033]$. The largest certified upper bound is $0.50050$. The second statement
follows from Lemmas~\ref{lem:mono} and~\ref{lem:bounds}.
\end{proof}

\begin{lemma}[middle region]\label{lem:middle}
$Q_\varphi(\alpha,x)>0$ for $\alpha\in[4,8.8033]$ and $e^{\psi(1+\alpha)}\le x\le16$.
\end{lemma}
\begin{proof}
Each of the first four ranges is cut into $20$ equal slabs and $[8,8.8033]$ into $16$ ($96$
slabs, of width at most $0.0503$). On a slab with left end
$a$, the $x$-interval runs from a certified lower bound of $e^{\psi(1+a)}$ to $16$. Boxes are
bisected adaptively. On a box the program bounds $T=e^{-x}S_\varphi$ from below by the larger of a
direct ball evaluation and the mean-value form
\[
 T(\alpha_c,x_c)-\sup|\partial_\alpha T|\,r_\alpha-\sup|\partial_xT|\,r_x ,
\]
where $(\alpha_c,x_c)$ is the centre, $r_\alpha$ and $r_x$ are the half-widths, and the suprema are
enclosures obtained by forward-mode differentiation in ball arithmetic. It also bounds
$\alpha^2e^{-x}B_\varphi$ from above over the box. A box is accepted when the resulting lower bound
of $Q_\varphi$ is positive. Otherwise the direction with the larger first-order term is bisected,
and a depth guard of $40$ can only cause failure. The run accepts $385$ boxes in $674$
evaluations, with maximum depth $18$. The smallest certified lower bound is $1.2101\cdot10^{-5}$, on
$\alpha\in[8.803104,8.8033]$, $x\in[10.75224,10.75883]$, next to the critical point.
\end{proof}

\begin{lemma}[witness]\label{lem:witness}
At $\alpha=8.8034$ and $x=10.7539$,
$\partial_\alpha E_\alpha(x^\alpha)\in[0.003968997501\pm3.5\cdot10^{-13}]$; in particular it is
positive.
\end{lemma}
\begin{proof}
Use the series of Lemma~\ref{lem:small} directly: the first $40$ terms are enclosed in Arb. For
$n\ge41$ the $n$-th term has modulus at most $n\bigl(\log(1+\alpha n)+|\log x|\bigr)(ex/(\alpha n))^{\alpha n}$,
by $\Gamma(1+y)\ge(y/e)^y$ and $0<\psi(y)\le\log y$. The ratio of consecutive bounds is below
$\frac12$, so the tail is at most twice its first term, which is below $6.9\cdot10^{-392}$. This
lemma does not use Lemma~\ref{lem:rays}.
\end{proof}

\begin{proof}[Proof of Theorem~\ref{thm:sharp}]
(i) On $(0,2)$ use \cite[Theorem~A]{FS}, whose proof gives strict inequalities, and on $[2,4]$
Theorem~\ref{thm:24}. Let
$\alpha\in[4,8.8033]$ and $x>0$. If $x\le e^{\psi(1+\alpha)}$, Lemma~\ref{lem:small} applies. If
$e^{\psi(1+\alpha)}\le x\le16$, Lemmas~\ref{lem:middle} and~\ref{lem:bounds} give
$\alpha^2e^{-x}H_\alpha(x)\ge Q_\varphi(\alpha,x)>0$. If $x\ge16$, use Lemma~\ref{lem:tail16}.
Hence $H_\alpha(x)>0$ on $[2,8.8033]\times(0,\infty)$, and continuity of $\alpha\mapsto U_\alpha(x)$
at $\alpha=2$ gives strict decrease on $(0,8.8033]$.

(ii) This is Lemma~\ref{lem:witness}. By continuity, $\partial_\alpha U_\alpha(10.7539)>0$ on a
neighbourhood of $8.8034$, so $U_\alpha(10.7539)$ increases on an interval to the left of
$8.8034$. Hence $\alpha^*<8.8034$; part~(i) gives $\alpha^*\ge8.8033$.
\end{proof}

\subsection{The critical point}\label{subsec:critical}
The following values are numerical (40-digit arithmetic) and are not used in the proofs. Solving
$\partial_\alpha U_\alpha(x)=0$, $\partial_x\partial_\alpha U_\alpha(x)=0$ near the witness gives
\[
 \alpha^*=8.80337769817047\ldots,\qquad x^*=10.7538996236407\ldots,
\]
with $\partial_x^2\partial_\alpha U=-287.96\ldots<0$ and $\partial_\alpha^2U=177.97\ldots>0$ at this
point. The first loss of monotonicity is therefore a nondegenerate fold at a moderate value of $x$.
In terms of~\eqref{eq:rays}, the residue sum from $w_{\pm1}$, $w_{\pm2}$ and $w_{\pm3}$ cancels the
main term $e^x$ exactly there. The natural guess that $\alpha=4$ is the boundary (where the
corrections stop being bounded in $x$) is thus not the true threshold. Sampling on a grid of step
$0.01$ in $\alpha$ shows a point with $\partial_\alpha U_\alpha(x)>0$ for every $\alpha\in[8.81,30]$;
this is not certified. The example $U_{25}(30)<U_{26}(30)$ of \cite{FS} lies in this region.

\section{Proofs of the corollaries}\label{sec:cor}
\begin{proof}[Proof of Corollary~\ref{cor:beta}]
For $\beta=1$ this is Theorem~\ref{thm:sharp}(i). For $\beta>1$, the positive series and
$B(1+\alpha n,\beta-1)=\Gamma(1+\alpha n)\Gamma(\beta-1)/\Gamma(\beta+\alpha n)$ give
\[
 E_{\alpha,\beta}(x^\alpha)=\frac{1}{\Gamma(\beta-1)}\int_0^1(1-t)^{\beta-2}E_\alpha\bigl((tx)^\alpha\bigr)\,dt .
\]
If $0<\alpha_1<\alpha_2\le8.8033$, the integrand for $\alpha_1$ is larger than that for $\alpha_2$
at every $t\in(0,1)$ by Theorem~\ref{thm:sharp}(i), and integration preserves the strict
inequality. This argument is taken from the draft, where it was used on $(0,4]$.
\end{proof}

\begin{proof}[Proof of Corollary~\ref{cor:bounds}]
By Theorem~\ref{thm:sharp}(i), $U_2>U_\alpha>U_4$ for $2<\alpha<4$, and $U_4>U_\alpha>U_8$ for
$4<\alpha<8$. For an integer $n\ge1$, summing the exponential series over residue classes modulo
$n$ gives $U_n(x)=\frac1n\sum_{k=0}^{n-1}\exp(xe^{2\pi ik/n})$. For $n=2,4,8$ this yields $\cosh x$,
$\frac12(\cosh x+\cos x)$ and
$\frac14\bigl(\cosh x+\cos x+2\cos\frac{x}{\sqrt2}\cosh\frac{x}{\sqrt2}\bigr)$.
\end{proof}

\section{Verification record}\label{sec:record}
Table~\ref{tab:record} lists what each program establishes. The programs print the numbers quoted
in Sections~\ref{sec:A}--\ref{sec:sharp}. No random samples enter any certified step.

\begin{table}[ht]\centering\small\setlength{\tabcolsep}{4pt}
\begin{tabular}{P{2.4cm}P{2.6cm}P{1.8cm}P{5.6cm}r}
\toprule
Claim & Program & Arithmetic & Output (abridged) & Time\\
\midrule
Lemmas~\ref{lem:scalar}, \ref{lem:certA} & Appendix~\ref{app:A} & Arb, 96 bits &
 $16{,}749$ rectangles, $33{,}497$ evaluations, depth $16$, area $27$, minimum $1.149\cdot10^{-5}$ & 2.6 s\\
Lemma~\ref{lem:certB} & Appendix~\ref{app:B} & Arb, 128 bits &
 maxima $0.199825$ and $0.182407$ & 0.1 s\\
Lemmas~\ref{lem:geom}--\ref{lem:witness} & Appendix~\ref{app:C} & Arb, 128 bits &
 $G(\cdot,16)\le0.50050$; $385$ boxes, $674$ evaluations, depth $18$, minimum $Q$ $1.2101\cdot10^{-5}$;
 witness $0.003968997501\pm3.5\cdot10^{-13}$ & 0.1 s\\
Cross-check of \eqref{eq:rays} (not part of the proof) & Appendices~\ref{app:B} and~\ref{app:C} with \texttt{-{}-numeric} & mpmath, 40 digits &
 $36$ and $32$ cases, relative differences $\le9.6\cdot10^{-41}$ and $\le3.5\cdot10^{-37}$ & about 2 min\\
\bottomrule
\end{tabular}
\caption{Verification record (Python~3.13, python-flint~0.9.0, mpmath~1.4.1).}\label{tab:record}
\end{table}

To reproduce, install \texttt{python-flint} and \texttt{mpmath}, save the three listings of
Appendix~\ref{app:code} under the names shown, and run
\begin{quote}\ttfamily\small
python3 verify\_bromwich\_certificate.py\\
python3 verify\_rays\_2\_4.py -{}-numeric\\
python3 verify\_sharp\_range.py -{}-numeric
\end{quote}
The software version matters because the certificate relies on special-function enclosures;
the versions used are recorded above.

\begin{thebibliography}{9}
\bibitem{FS}
R.~A.~C. Ferreira and T.~Simon, \emph{Mittag--Leffler functions and convex ordering},
arXiv:2512.04940v1, 4 December 2025.
\href{https://arxiv.org/abs/2512.04940v1}{arxiv.org/abs/2512.04940v1}.
Theorem~A and Remark~1(a).

\bibitem{GKMR}
R.~Gorenflo, A.~A. Kilbas, F.~Mainardi and S.~V. Rogosin,
\emph{Mittag--Leffler Functions, Related Topics and Applications}, 2nd ed.,
Springer Monographs in Mathematics, Springer, Berlin, 2020.

\bibitem{Arb}
F.~Johansson, \emph{Arb: efficient arbitrary-precision midpoint-radius interval arithmetic},
IEEE Transactions on Computers 66 (2017), 1281--1292.
\href{https://doi.org/10.1109/TC.2017.2690633}{doi:10.1109/TC.2017.2690633}.

\bibitem{flint}
The python-flint developers, \emph{python-flint 0.9.0}, Python interface to FLINT and Arb.
\href{https://python-flint.readthedocs.io/}{python-flint.readthedocs.io}.
\end{thebibliography}

\appendix
\section{Programs}\label{app:code}
\subsection{Proof A: compact certificate (\texttt{verify\_bromwich\_certificate.py})}\label{app:A}
\begin{lstlisting}
#!/usr/bin/env python3
"""
Independent re-implementation of the compact certificate of Proof A (Bromwich-line proof),
Lemmas 3.2-3.3 of the paper:

    for every (alpha, x) in [2,4] x [5/2,16]:   e^{-x} H_{alpha,24}(x) > 1/100000,
    H_{alpha,24}(x) = sum_{n=1}^{24} n x^{alpha n} (psi(1 + alpha n) - log x) / Gamma(1 + alpha n),

together with the two scalar inequalities psi(3) > log(5/2) and psi(51) > log(16).

Arb ball arithmetic (python-flint), exact rational box endpoints, adaptive bisection
(alpha is bisected when 8 * width(alpha) > width(x), otherwise x), depth guard = failure.
Coverage invariant: every popped box is either accepted or replaced by two halves.
"""
import sys
import time
from flint import arb, fmpq, ctx

ctx.prec = 96
THRESHOLD = arb(fmpq(1, 100000))


def ball(lo, hi):
    return arb(lo).union(arb(hi))


def lower_bound(alo, ahi, xlo, xhi):
    A = ball(alo, ahi)
    X = ball(xlo, xhi)
    lx = X.log()
    s = arb(0)
    for n in range(1, 25):
        nA = n * A
        s += n * (nA * lx - (1 + nA).lgamma() - X).exp() * ((1 + nA).digamma() - lx)
    return s.lower()


def scalar_inequalities():
    a = arb(3).digamma() - arb(fmpq(5, 2)).log()
    b = arb(51).digamma() - arb(16).log()
    return a, b, (a > 0) and (b > 0)


def certify():
    stack = [(fmpq(2), fmpq(4), fmpq(5, 2), fmpq(16), 0)]
    accepted = evals = maxdepth = 0
    area = fmpq(0)
    minlb = None
    while stack:
        a0, a1, x0, x1, d = stack.pop()
        if d > 60:
            return None
        evals += 1
        lb = lower_bound(a0, a1, x0, x1)
        if lb > THRESHOLD:
            accepted += 1
            maxdepth = max(maxdepth, d)
            area += (a1 - a0) * (x1 - x0)
            if minlb is None or lb < minlb:
                minlb = lb
            continue
        if 8 * (a1 - a0) > (x1 - x0):
            am = (a0 + a1) / 2
            stack += [(a0, am, x0, x1, d + 1), (am, a1, x0, x1, d + 1)]
        else:
            xm = (x0 + x1) / 2
            stack += [(a0, a1, x0, xm, d + 1), (a0, a1, xm, x1, d + 1)]
    return accepted, evals, maxdepth, area, minlb


if __name__ == '__main__':
    t0 = time.time()
    a, b, ok = scalar_inequalities()
    print(f'psi(3) - log(5/2) = {a.str(6)},  psi(51) - log 16 = {b.str(6)}  -> {"OK" if ok else "FAIL"}')
    res = certify()
    if res is None or not ok:
        print('CERTIFICATE FAILED')
        sys.exit(1)
    acc, ev, md, area, mlb = res
    print(f'accepted rectangles = {acc}, evaluations = {ev}, max depth = {md}, '
          f'exact covered area = {area} (target 27), full binary tree: {ev == 2 * acc - 1}')
    print(f'smallest certified lower bound of e^-x H_24 = {mlb.str(5, radius=False)}  (> 1/100000)')
    print(f'elapsed {time.time() - t0:.1f} s')
\end{lstlisting}
\subsection{Proof B: one-dimensional certificate (\texttt{verify\_rays\_2\_4.py})}\label{app:B}
\begin{lstlisting}
#!/usr/bin/env python3
"""
Proof B of Theorem 1.1 (rotated rays), computer-assisted step = Lemma 4.1 of the paper:
for every x > 0, alpha -> E_alpha(x^alpha) is strictly decreasing on [2, 4]
(Ferreira-Simon, arXiv:2512.04940, Remark 1(a)).

Usage
    pip install python-flint mpmath
    python3 verify_rays_2_4.py            # rigorous check (Lemma 4.1), ~0.1 s
    python3 verify_rays_2_4.py --numeric  # also: ray formula (Lemma 2.5) vs series, ~5 s

Notation (all logs/powers principal):  theta = 2*pi/alpha,  lambda = -x*cos(phi),  gamma = alpha*phi,
    m2 = 1 if cos(gamma) <= 0 else sin(gamma)^2,
    B  = [ Gamma(a) lambda^-a (phi + (psi(a) - ln lambda)_+) + 2 min(1/a^2, Gamma(a-1)/(e lambda^(a-1))) ] / (pi m2),
    G  = 1[theta < phi] * 2 exp(-x(1 - cos theta)) (1 + theta x) + a^2 e^-x B.
Lemma 2.6: D(a,x) := d/da E_a(x^a) <= (e^x / a^2) (G - 1).   Lemma 2.7: G is decreasing in x for x >= 1.
Lemma 4.1 (this script): on each sub-interval [a, a'] of [2, 2.5] (phi = 0.7 pi) and [2.5, 4] (phi = 0.9 pi),
    sup_{alpha in [a,a']} G(alpha, X) < 1,  where X is a rigorous lower bound of exp(psi(1 + a)).
All quantities are enclosed with Arb ball arithmetic (python-flint), so rounding errors are accounted for.
"""
import sys
from flint import arb, ctx

ctx.prec = 128
PI = arb.pi()
E = arb(1).exp()


def G_upper(a_lo, a_hi, phi_frac, residue_inside):
    """Rigorous upper bound of sup_{alpha in [a_lo, a_hi]} G(alpha, X), X <= exp(psi(1 + a_lo))."""
    a_lo, a_hi = arb(a_lo), arb(a_hi)
    alpha = arb(((a_lo + a_hi) / 2).mid(), ((a_hi - a_lo) / 2).upper())  # ball containing [a_lo, a_hi]
    assert alpha.contains(a_lo) and alpha.contains(a_hi)
    phi = PI * arb(phi_frac)
    X = (1 + a_lo).digamma().exp().lower()          # X <= e^{psi(1+alpha)} for every alpha in the interval
    assert X > 1                                    # needed for Lemma 2.7
    theta = 2 * PI / alpha

    term_res = arb(0)
    if residue_inside:                              # poles e^{+-i theta} lie inside |arg w| < phi
        assert theta < phi
        term_res = 2 * (-(X * (1 - theta.cos()))).exp() * (1 + theta * X)
    else:                                           # poles lie outside: no residue term
        assert theta > phi

    gam = alpha * phi
    if gam.cos().upper() <= 0:
        m2 = arb(1)
    else:
        m2 = (gam.sin() ** 2).lower()
        assert m2 > 0
    lam = -X * phi.cos()
    assert lam > 0
    lnlam = lam.log()
    p1 = (alpha.gamma().upper() * (-(alpha * lnlam)).exp().upper()
          * (phi.upper() + max((alpha.digamma() - lnlam).upper(), arb(0))))
    lb = min((1 / alpha ** 2).upper(), ((alpha - 1).gamma() / (E * lam ** (alpha - 1))).upper())
    B = (p1 + 2 * lb) / (PI.lower() * m2)
    term_ray = (alpha ** 2).upper() * (-X).exp().upper() * B
    return (term_res.upper() + term_ray.upper()).upper()


def rigorous_check():
    ok = True
    for lo, hi, n, phi_frac, inside, name in [(2, '2.5', 500, '0.7', False, 'Range I  [2,2.5], phi=0.7pi'),
                                               ('2.5', 4, 1500, '0.9', True, 'Range II [2.5,4], phi=0.9pi')]:
        worst, where = arb(0), None
        lo_, hi_ = arb(lo), arb(hi)
        for i in range(n):
            a0 = (lo_ + (hi_ - lo_) * i / n).mid()
            a1 = (lo_ + (hi_ - lo_) * (i + 1) / n).mid()
            g = G_upper(a0, a1, phi_frac, inside)
            if not g < 1:
                ok = False
                print('FAIL', a0, a1, g)
            if g > worst:
                worst, where = g, (a0, a1)
        print(f'{name}: {n} sub-intervals, max upper bound of G = {worst.str(6, radius=False)} '
              f'on [{where[0].str(6, radius=False)}, {where[1].str(6, radius=False)}]')
    print('Lemma 4.1 VERIFIED: G < 1 on all of [2,4]' if ok else 'Lemma 4.1 FAILED')
    return ok


def numeric_contour_check():
    """Non-rigorous sanity check of Lemma 2.5: ray formula vs. term-wise differentiated series."""
    import mpmath as mp
    mp.mp.dps = 40

    def D_series(a, x):
        a, x = mp.mpf(a), mp.mpf(x)
        lx, s, n = mp.log(x), mp.mpf(0), 1
        while True:
            an = a * n
            t = mp.exp(an * lx - mp.loggamma(1 + an))
            s += n * (lx - mp.digamma(1 + an)) * t
            if an > x + 60 and t < mp.mpf(10) ** -50:
                return s
            n += 1

    def D_contour(a, x, phi):
        a, x, phi = mp.mpf(a), mp.mpf(x), mp.mpf(phi)
        th = 2 * mp.pi / a
        res = mp.exp(x) / a ** 2
        if th < phi:
            w = mp.exp(1j * th)
            res += 2 * mp.re(mp.exp(x * w) / a ** 2 * (1 + x * w * 1j * th))
        e, eap = mp.exp(1j * phi), mp.exp(1j * a * phi)
        f = lambda r: mp.exp(x * r * e) * r ** (a - 1) * eap * (mp.log(r) + 1j * phi) / (r ** a * eap - 1) ** 2
        J = mp.im(mp.quad(f, [0, 0.5, 1, 2, 5, mp.inf])) / mp.pi
        return -res - J

    worst = mp.mpf(0)
    for a in ['2.1', '2.3', '2.7', '3.0', '3.4', '3.9']:
        for x in ['0.7', '2.5', '6']:
            ds = D_series(a, x)
            for phi in [0.75 * mp.pi, 0.92 * mp.pi]:
                worst = max(worst, abs(ds - D_contour(a, x, phi)) / abs(ds))
    print(f'Lemma 2.5 numeric check: 36 (alpha, x, phi) cases, max relative difference = {mp.nstr(worst, 3)}')


if __name__ == '__main__':
    ok = rigorous_check()
    if '--numeric' in sys.argv:
        numeric_contour_check()
    sys.exit(0 if ok else 1)
\end{lstlisting}
\subsection{Sharp range (\texttt{verify\_sharp\_range.py})}\label{app:C}
\begin{lstlisting}
#!/usr/bin/env python3
"""
Theorem 1.2 (sharp range), computer-assisted steps = Lemmas 5.1-5.4 of the paper.
Rigorous (Arb ball arithmetic via python-flint) certificate that
    D(alpha, x) := d/dalpha E_alpha(x^alpha) < 0   for all alpha in [4, 88033/10000], x > 0,
plus a certified witness D(8.8034, 10.7539) > 0.  Hence 8.8033 <= alpha* < 8.8034.
Usage:  python3 verify_sharp_range.py [--numeric]   (~0.1 s; --numeric adds a ~2 min series cross-check)

Representation used for x >= e^{psi(1+alpha)} (Lemma 2.5; cross-checked against the series by --numeric):
    D = -(e^x + S(alpha, x)) / alpha^2 - J_phi(alpha, x),
    S = sum_{k=1}^{K} 2 e^{x cos u_k} [cos(x sin u_k) - u_k x sin(x sin u_k + u_k)],  u_k = 2 pi k / alpha,
    |J_phi| <= B(alpha, x)  (Lemma 2.6; m^2 > 1/2 asserted on every range below = Lemma 5.1).
Hence D < 0 whenever Q := 1 + e^{-x} S - alpha^2 e^{-x} B > 0.

Regions:  x <= e^{psi(1+alpha)}: termwise sign (no computation).
          e^{psi(1+alpha)} <= x <= 16: 2-D adaptive certificate of Q > 0 (Lemma 5.3).
          x >= 16: analytic bound G = sum_k 2 e^{-x(1-cos u_k)}(1 + u_k x) + alpha^2 e^{-x} B < 1,
                   decreasing in x for x >= 16 (Lemma 2.7); certified at x = 16 on sub-intervals (Lemma 5.2).
Witness:  series partial sum + explicit tail bound, independent of Lemma 2.5 (Lemma 5.4).
"""
import sys
import time
from flint import arb, fmpq, ctx

ctx.prec = 128
PI = arb.pi()
E = arb(1).exp()
X_TOP = arb(16)

# (alpha_lo, alpha_hi, phi / pi, K = number of residue pairs inside |arg w| < phi)
RANGES = [
    (fmpq(4), fmpq(5), fmpq(13, 20), 1),
    (fmpq(5), fmpq(6), fmpq(9, 10), 2),
    (fmpq(6), fmpq(7), fmpq(19, 25), 2),
    (fmpq(7), fmpq(8), fmpq(33, 50), 2),
    (fmpq(8), fmpq(88033, 10000), fmpq(83, 100), 3),
]
WITNESS = (fmpq(44017, 5000), fmpq(107539, 10000))   # (alpha, x) = (8.8034, 10.7539)


def ball(lo, hi):
    """Ball containing the closed interval [lo, hi] (fmpq or arb endpoints)."""
    a = lo if isinstance(lo, arb) else arb(lo)
    b = hi if isinstance(hi, arb) else arb(hi)
    return a.union(b)


class Dual:
    """Value with enclosures of the partial derivatives d/dalpha and d/dx."""
    __slots__ = ('v', 'a', 'x')

    def __init__(self, v, a=None, x=None):
        self.v = v
        self.a = a if a is not None else arb(0)
        self.x = x if x is not None else arb(0)

    def __add__(self, o):
        o = lift(o)
        return Dual(self.v + o.v, self.a + o.a, self.x + o.x)

    __radd__ = __add__

    def __sub__(self, o):
        o = lift(o)
        return Dual(self.v - o.v, self.a - o.a, self.x - o.x)

    def __rsub__(self, o):
        return lift(o) - self

    def __mul__(self, o):
        o = lift(o)
        return Dual(self.v * o.v, self.a * o.v + self.v * o.a, self.x * o.v + self.v * o.x)

    __rmul__ = __mul__

    def __neg__(self):
        return Dual(-self.v, -self.a, -self.x)


def lift(o):
    if isinstance(o, Dual):
        return o
    return Dual(o if isinstance(o, arb) else arb(o))


def dexp(d):
    e = d.v.exp()
    return Dual(e, e * d.a, e * d.x)


def dsin(d):
    c = d.v.cos()
    return Dual(d.v.sin(), c * d.a, c * d.x)


def dcos(d):
    s = d.v.sin()
    return Dual(d.v.cos(), -s * d.a, -s * d.x)


def dinv(d):
    i = 1 / d.v
    return Dual(i, -(i * i) * d.a, -(i * i) * d.x)


def T_dual(A, X, K):
    """e^{-x} S(alpha, x) with derivative enclosures, for alpha in A, x in X."""
    al = Dual(A, arb(1), arb(0))
    xx = Dual(X, arb(0), arb(1))
    th = dinv(al) * (2 * PI)
    tot = Dual(arb(0))
    for k in range(1, K + 1):
        u = th * k
        c = dcos(u)
        s = dsin(u)
        e = dexp(-(xx * (1 - c)))
        p = xx * s
        tot = tot + e * (dcos(p) - u * xx * dsin(p + u)) * 2
    return tot


def m2_lower(A, phi):
    """Lower bound of m^2 = min_r |r^alpha e^{i alpha phi} - 1|^2: 1 if cos(alpha phi) <= 0, else sin^2(alpha phi)."""
    g = A * phi
    if g.cos().upper() <= 0:
        return arb(1)
    m2 = min(arb(1), (g.sin() ** 2).lower())
    assert m2 > 0
    return m2


def B_upper(A, X, phi):
    """Upper bound of alpha^2 e^{-x} B(alpha, x) over the box."""
    m2 = m2_lower(A, phi)
    lam = -X * phi.cos()
    assert lam > 0
    lnlam = lam.log()
    p1 = A.gamma() * (-(A * lnlam)).exp() * (phi + max((A.digamma() - lnlam).upper(), arb(0)))
    l1 = (1 / A ** 2).upper()
    l2 = ((A - 1).gamma() / (E * (lam ** (A - 1)))).upper()
    B = (p1 + 2 * min(l1, l2)) / (PI.lower() * m2)
    return ((A ** 2) * (-X).exp() * B).upper()


def check_range_geometry(lo, hi, phi_frac, K):
    A = ball(lo, hi)
    phi = PI * arb(phi_frac)
    th = 2 * PI / A
    assert phi > PI / 2 and phi < PI
    for k in range(1, K + 1):
        assert k * th < phi, ('pole should be inside', k)
    assert (K + 1) * th > phi, 'next pole must be outside the sector'
    n = 400
    for i in range(n):
        Ai = ball(lo + (hi - lo) * i / n, lo + (hi - lo) * (i + 1) / n)
        assert m2_lower(Ai, phi) > arb(1) / 2, 'rays must stay well away from the poles'
    return phi


def Q_lower(alo, ahi, xlo, xhi, phi, K):
    A = ball(alo, ahi)
    X = ball(xlo, xhi)
    tb = T_dual(A, X, K)
    naive = tb.v.lower()
    ac = arb((alo + ahi) / 2)
    xc = (xlo + xhi) / 2
    ra = arb((ahi - alo) / 2)
    rx = (xhi - xlo) / 2
    tc = T_dual(ac, xc, K).v
    centered = (tc - tb.a.abs_upper() * ra - tb.x.abs_upper() * rx).lower()
    t_low = max(naive, centered)
    q = (1 + t_low - B_upper(A, X, phi)).lower()
    return q, tb.a.abs_upper() * ra, tb.x.abs_upper() * rx


def certify_middle():
    stats = dict(accepted=0, evals=0, maxdepth=0, minq=arb(10), minq_box=None)
    for lo, hi, phi_frac, K in RANGES:
        phi = check_range_geometry(lo, hi, phi_frac, K)
        n_slabs = int(((hi - lo) * 20).floor()) or 1           # slabs of width 1/20
        for i in range(n_slabs):
            alo = lo + (hi - lo) * i / n_slabs
            ahi = lo + (hi - lo) * (i + 1) / n_slabs
            xlo = (1 + arb(alo)).digamma().exp().lower()       # <= e^{psi(1+alpha)} on the slab
            stack = [(alo, ahi, xlo, X_TOP, 0)]
            while stack:
                a0, a1, x0, x1, depth = stack.pop()
                if depth > 40:
                    print('FAIL: depth guard at', a0, a1, x0, x1)
                    return None
                stats['evals'] += 1
                q, ea, ex = Q_lower(a0, a1, x0, x1, phi, K)
                if q > 0:
                    stats['accepted'] += 1
                    stats['maxdepth'] = max(stats['maxdepth'], depth)
                    if q < stats['minq']:
                        stats['minq'] = q
                        stats['minq_box'] = (a0, a1, x0, x1)
                    continue
                if ea > ex:
                    am = (a0 + a1) / 2
                    stack += [(a0, am, x0, x1, depth + 1), (am, a1, x0, x1, depth + 1)]
                else:
                    xm = (x0 + x1) / 2
                    stack += [(a0, a1, x0, xm, depth + 1), (a0, a1, xm, x1, depth + 1)]
    return stats


def certify_tail():
    worst = arb(0)
    for lo, hi, phi_frac, K in RANGES:
        phi = check_range_geometry(lo, hi, phi_frac, K)
        n = int(((hi - lo) * 200).floor()) or 1
        for i in range(n):
            A = ball(lo + (hi - lo) * i / n, lo + (hi - lo) * (i + 1) / n)
            th = 2 * PI / A
            g = arb(0)
            for k in range(1, K + 1):
                u = k * th
                g += 2 * (-(X_TOP * (1 - u.cos()))).exp() * (1 + u * X_TOP)
            G = (g.upper() + B_upper(A, X_TOP, phi)).upper()
            assert G < 1, ('tail bound fails', i)
            worst = max(worst, G)
    return worst


def numeric_representation_check():
    """Non-rigorous sanity check of the residue + ray representation against the series (mpmath, 40 digits)."""
    import mpmath as mp
    mp.mp.dps = 40

    def D_series(a, x):
        a, x = mp.mpf(a), mp.mpf(x)
        lx, s, n = mp.log(x), mp.mpf(0), 1
        while True:
            an = a * n
            t = mp.exp(an * lx - mp.loggamma(1 + an))
            s += n * (lx - mp.digamma(1 + an)) * t
            if an > x + 60 and t < mp.mpf(10) ** -50:
                return s
            n += 1

    def D_rep(a, x, phi):
        a, x, phi = mp.mpf(a), mp.mpf(x), mp.mpf(phi)
        th, S, k = 2 * mp.pi / a, mp.mpf(0), 1
        while k * th < phi:
            u = k * th
            S += 2 * mp.exp(x * mp.cos(u)) * (mp.cos(x * mp.sin(u)) - u * x * mp.sin(x * mp.sin(u) + u))
            k += 1
        e, eap = mp.exp(1j * phi), mp.exp(1j * a * phi)
        f = lambda r: mp.exp(x * r * e) * r ** (a - 1) * eap * (mp.log(r) + 1j * phi) / (r ** a * eap - 1) ** 2
        J = mp.im(mp.quad(f, [0, 0.5, 1, 2, 5, 20, mp.inf])) / mp.pi
        return -(mp.exp(x) + S) / a ** 2 - J

    worst = mp.mpf(0)
    for a, pf in [(4.5, 0.65), (5.5, 0.9), (6.0, 0.76), (6.5, 0.76), (7.5, 0.66), (8.0, 0.83), (8.5, 0.83), (8.8, 0.83)]:
        for x in [1.5, 6.0, 10.75, 14.0]:
            ds = D_series(a, x)
            worst = max(worst, abs(ds - D_rep(a, x, pf * mp.pi)) / abs(ds))
    print(f'representation check (32 cases, alpha in [4.5, 8.8]): max relative difference = {mp.nstr(worst, 3)}')


def witness(alpha, x, N=40):
    """Rigorous enclosure of D(alpha, x) from the series: partial sum to N plus tail bound."""
    A = arb(alpha)
    X = arb(x)
    lx = X.log()
    s = arb(0)
    for n in range(1, N + 1):
        an = A * n
        s += n * (lx - (1 + an).digamma()) * (an * lx - (1 + an).lgamma()).exp()
    # tail: |term_n| <= n (log(1+an) + |log x|) (e x / (a n))^(a n)  (Stirling: Gamma(1+y) >= (y/e)^y)
    bnd = []
    for n in (N + 1, N + 2):
        an = A * n
        bnd.append(n * ((1 + an).log() + lx.abs_upper()) * ((E * X / an) ** an))
    ratio = (bnd[1] / bnd[0]).upper()
    assert ratio < arb(1) / 2
    tail = (2 * bnd[0]).upper()
    return s, tail


if __name__ == '__main__':
    # optional:  --numeric (representation vs series, ~2 min)  --amax p/q  (top of the certified alpha range)   --witness p/q,r/s  (alpha, x) of the failure witness
    args = sys.argv[1:]
    if '--amax' in args:
        p, q = args[args.index('--amax') + 1].split('/')
        lo, _, ph, K = RANGES[-1]
        RANGES[-1] = (lo, fmpq(int(p), int(q)), ph, K)
    if '--witness' in args:
        wa, wx = args[args.index('--witness') + 1].split(',')
        WITNESS = tuple(fmpq(*map(int, w.split('/'))) for w in (wa, wx))
    t0 = time.time()
    tail_worst = certify_tail()
    print(f'certified alpha range: [4, {float(RANGES[-1][1])}]')
    print(f'x >= 16: analytic bound G < 1 on the whole range; max certified G(alpha,16) = {tail_worst.str(5, radius=False)}')
    st = certify_middle()
    if st is None:
        print('MIDDLE REGION FAILED')
        sys.exit(1)
    b = st['minq_box']
    print(f"middle region e^psi(1+a) <= x <= 16: accepted boxes = {st['accepted']}, evaluations = {st['evals']}, "
          f"max depth = {st['maxdepth']}")
    print(f"   smallest certified lower bound of Q = {st['minq'].str(5, radius=False)} on alpha in "
          f"[{float(b[0]):.6f}, {float(b[1]):.6f}], x in [{float(b[2].mid()):.5f}, {float(b[3].mid()):.5f}]")
    s, tail = witness(*WITNESS)
    print(f'witness D({float(WITNESS[0])}, {float(WITNESS[1])}): partial sum = {s.str(10)}, tail <= {tail.str(3, radius=False)}; '
          f'certified positive: {(s - tail) > 0}')
    print(f'elapsed {time.time() - t0:.1f} s')
    if '--numeric' in args:
        numeric_representation_check()
\end{lstlisting}


\end{document}
